\documentclass[11pt]{article}
\usepackage[margin=1in]{geometry}
\usepackage{amsmath, amssymb}
\usepackage{graphicx}
\usepackage{array}
\usepackage{tabularx}
\usepackage{booktabs}
\usepackage{enumitem}
\usepackage{hyperref}
\usepackage{caption}
\usepackage{titlesec}
\usepackage{float}
\titleformat{\section}{\large\bfseries}{\thesection.}{1em}{}

\begin{document}

\begin{center}
    \Large \textbf{Sri Sivasubramaniya Nadar College of Engineering, Chennai} \\
    \large (An Autonomous Institution affiliated to Anna University) \\
    \vspace{0.3cm}
\end{center}

\begin{table}[h!]
\renewcommand{\arraystretch}{1.5}
\centering
\resizebox{\textwidth}{!}{%
\begin{tabular}{|l|cll|}
\hline
\textbf{Degree \& Branch} & \multicolumn{1}{c|}{B.E Computer Science \& Engineering} & \textbf{Semester} & VI \\ \hline
\textbf{Subject Code \& Name} & \multicolumn{3}{c|}{UCS2612 -- Machine Learning Laboratory} \\ \hline
\textbf{Academic Year} & \multicolumn{1}{c|}{2025--2026 (Even)} & \textbf{Batch} & 2023--2027 \\ \hline
\textbf{Name} & \multicolumn{1}{c|}{Mehanth T} & \textbf{Register No.} & 3122235001080 \\ \hline
\end{tabular}
}
\end{table}

\begin{center}
    \textbf{\Large Experiment 8}
\end{center}

\begin{center}
    \textbf{Effect of PCA on Regression and Classification Models}
\end{center}

\section*{Objective}
To study the impact of \textbf{Principal Component Analysis (PCA)} on the performance of machine learning models.
\begin{enumerate}[label=\alph*)]
    \item Train models in the original feature space (Without PCA)
    \item Train models in reduced feature space (With PCA)
    \item Compare performance using 5-fold cross-validation
\end{enumerate}

\section*{Dataset}
\textbf{Dataset: PCA\_Full\_Lab\_Dataset\_1000\_Samples.csv}

\begin{itemize}
    \item Total samples: 1000
    \item Total features: 11 (Math, Physics, Chemistry, Biology, Total\_Science, Average\_Science, English, Computer\_Science, Sports\_Score, Attendance\_Percentage, Random\_Noise\_Feature)
    \item \textbf{Regression Target:} Final\_Score\_Regression (continuous)
    \item \textbf{Classification Target:} Performance\_Level\_Classification (High / Medium / Low)
    \item Class distribution: High=202, Low=253, Medium=545
    \item Missing values: None
\end{itemize}

\section*{Preprocessing Steps}
\begin{enumerate}
    \item No missing values found; no imputation required.
    \item Encode classification target with LabelEncoder: Low=0, High=1, Medium=2.
    \item Standardise all features using \texttt{StandardScaler} (zero mean, unit variance) -- \emph{critical before PCA}.
    \item Evaluate using 5-fold cross-validation (KFold for regression, StratifiedKFold for classification).
\end{enumerate}

\section*{PCA Implementation}

\begin{enumerate}
    \item Apply PCA after standardisation.
    \item Retain 95\% of explained variance (automatic component selection).
    \item Plot Scree Plot showing cumulative variance vs.\ number of components.
\end{enumerate}

\begin{table}[H]
\centering
\caption{PCA Summary}
\begin{tabular}{|c|c|c|}
\hline
Components Chosen & Explained Variance (\%) & Justification \\
\hline
7 out of 11 & 95.90\% & Retains 95\% variance threshold; removes 4 redundant/noisy dimensions \\
\hline
\end{tabular}
\end{table}

\section*{Regression Results}

\begin{table}[H]
\centering
\caption{5-Fold CV Results -- Regression (1000 samples)}
\begin{tabularx}{\textwidth}{|X|X|X|X|X|}
\hline
\textbf{Model} & \textbf{Metric} & \textbf{Fold Avg (No-PCA)} & \textbf{Fold Avg (With-PCA)} & \textbf{Observation} \\
\hline
Linear Regression & MSE & 26.2338 & 26.2251 & Negligible improvement \\
Linear Regression & $R^2$ & 0.7479 & 0.7481 & Nearly identical \\
Random Forest & MSE & 29.0245 & 28.9563 & Slight improvement \\
Random Forest & $R^2$ & 0.7215 & 0.7222 & Minimal PCA effect \\
\hline
\end{tabularx}
\end{table}

\section*{Classification Results}

\begin{table}[H]
\centering
\caption{5-Fold CV Results -- Classification (1000 samples)}
\begin{tabularx}{\textwidth}{|X|X|X|X|X|}
\hline
\textbf{Model} & \textbf{Metric} & \textbf{Fold Avg (No-PCA)} & \textbf{Fold Avg (With-PCA)} & \textbf{Observation} \\
\hline
Logistic Regression & Accuracy & 0.7460 & 0.7460 & No change \\
Logistic Regression & F1-score  & 0.7437 & 0.7438 & Identical \\
SVM (RBF) & Accuracy & 0.7390 & 0.7340 & Slight decrease \\
SVM (RBF) & F1-score  & 0.7324 & 0.7264 & Slight decrease \\
\hline
\end{tabularx}
\end{table}

\section*{Analysis Questions}
\begin{enumerate}
    \item \textbf{Did PCA improve Linear Regression?} Negligibly (MSE 26.2338 $\to$ 26.2251). Since features were already clean and correlated, PCA had little effect on a linear model.
    \item \textbf{Did PCA significantly affect Random Forest?} No. Tree-based models handle high-dimensional and correlated features natively by splitting on the most informative features, so PCA provides minimal benefit.
    \item \textbf{Which classification model benefited more from PCA?} Neither model benefited. Logistic Regression was unchanged; SVM slightly decreased (0.7390 $\to$ 0.7340), possibly because PCA's rotation changed the RBF kernel's distance calculations.
    \item \textbf{Did PCA reduce variance across folds?} Marginally for regression (slight MSE reduction), but not significantly for classification in this dataset.
    \item \textbf{Why do linear models benefit more from PCA?} Linear models struggle with multicollinearity; PCA decorrelates features and removes noise dimensions. However, with only 11 features and moderate correlation, the gain was minimal here.
\end{enumerate}

\section*{Expected Learning Outcome}
\begin{itemize}
    \item PCA reduces multicollinearity and noise dimensions (11 $\to$ 7 components).
    \item PCA helps linear models more than tree-based models.
    \item For this dataset, features were already informative; PCA offered negligible benefit.
    \item Dimensionality reduction is most valuable with many correlated or noisy features.
    \item Always standardise before PCA.
\end{itemize}

\section*{Final Conclusion}
PCA compressed 11 features to 7 components retaining 95.90\% variance. For this dataset, PCA had negligible impact on all four models because the original features were already well-conditioned. PCA is most useful when: (1) features are highly correlated, (2) there are many noisy dimensions, or (3) computational cost reduction is critical. It is least necessary for tree-based models (Random Forest, Gradient Boosting) which handle correlated features inherently. The trade-off: PCA reduces dimensionality at the cost of interpretability -- components are linear combinations with no direct real-world meaning.

\end{document}
